\writeSectionFrame{\presentationTitle}{Fassade}
\begin{frame}{Steckbrief}
	\begin{itemize}
 		\item Deutscher Name: Fassade
 		\item Englischer Name: Facade
 		\item Gruppe: Strukturmuster
 	\end{itemize}
\end{frame}

\begin{frame}{Fassade}
	\begin{block}{Zweck}
 		Biete eine einheitliche Schnittstelle zu einer Menge von Schnittstellen eines Subsystems. Die Fassadenklasse definiert eine abstrakte Schnittstelle, welche die Verwendung des Subsystems vereinfacht.
 	\end{block}
\end{frame}

\frameWithImageOnly{\BILDERPFAD/fassade1.jpg}

\begin{frame}{Beispiel aus dem realen Leben}
	\begin{center}
 		\includegraphics[scale=0.25]{\BILDERPFAD/fassade2.png}
 	\end{center}
\end{frame}

\note{
	Die Fassade kann man mit einem Call-Center vergleichen. Wenn ich telefonisch eine 
	Bestellung aufgebe, lande ich in einem Call-Center. Irgend jemand aus dieser
	Fassade nimmt die Bestellung entgegen. Im Hintergrund wird dann ein riesiges
	Subsystem gestartet, aber davon bekomme ich als Kunde nichts mit. Das Call-Center ist 
	eine Fassade, die das ganze System vor mir versteckt. 
}

\begin{frame}{Aufbau von Fassaden}
	\begin{center}
 		\includegraphics[scale=0.25]{\BILDERPFAD/fassade3.png}
 	\end{center}
\end{frame}

\begin{frame}{Klassendiagramm Beispiel Versandhäuser}
    \begin{tikzpicture}[scale=0.7, transform shape]
        \umlclass[x=-3,y=0]{VersandhandelFassade}{
        }{
            +starteBestellung(bestellNummer: String)
        }

        \umlclass[x=4,y=0]{VersandhandelClient}{
        }{}

        \umlclass[x=-6,y=3]{Bestellannahme}{
            -bestellNummer: int
        }{
            
        }
        
        \umlclass[x=-6,y=-3]{Bestelleingang}{}{
            -bestellNummer: String
        }{
            +starteBestellung()
        }

        \umluniassoc{VersandhandelClient}{VersandhandelFassade}
        \umluniassoc{VersandhandelFassade}{Bestellannahme}
        \umluniassoc{VersandhandelFassade}{Bestelleingang}
    \end{tikzpicture}
\end{frame}
\note{
	Die System von Schleckermann und BestellFix sollen jetzt miteinander verbunden werden. Die Bestellnummern sind unterschiedlich aufgebaut, die von Schleckermann sind einfache Ganzzahlen, während die von BestellFix mit BF beginnen, und auch die Namen der Untermodule sind unterschiedlich. Die Fassade soll eine Bestellung anhand der Bestellnummer einem Sub-System zuordnen und die entsprechende Methode aufrufen, um eine Bestellung
	zu starten.
}
 
\begin{frame}[fragile]{Bestellklasse von Schleckermann}
	\begin{exampleblock}{Bestellannahme}
 		\begin{lstlisting}
public class Bestellannahme {
	private int bestellnummer;
	
	public Bestellannahme(int bestellnummer) {
		System.out.println("Bestellung mit der Bestellnummer " 
				+ bestellnummer + " wird bei Schleckermann gestartet.");
		this.bestellnummer = bestellnummer;
	}
	
	public int getBestellnummer() {
		return bestellnummer;
	}
}
 		\end{lstlisting}
 	\end{exampleblock}
\end{frame}

\begin{frame}[fragile]{Bestellklasse von BestellFix}
	\begin{exampleblock}{Bestelleingang}
 		\begin{lstlisting}
public class Bestelleingang {
	private String bestellNr;
	
	public void setBestellNr(String bestellNr) {
		this.bestellNr = bestellNr;
	}
	
	public void starteBestellung() {
		System.out.println("Bestellung bei BestellFix wird gestartet: " 
				+ bestellNr);
	}
}
 		\end{lstlisting}
 	\end{exampleblock}
\end{frame}

\begin{frame}[fragile]{Fassade}
	\begin{exampleblock}{VersandhandelFassade}
 		\begin{lstlisting}
public class VersandhandelFassade {
	private Bestellannahme bestellannahme;
	private Bestelleingang bestelleingang;
	
	public VersandhandelFassade() {
		bestelleingang = new Bestelleingang();
	}
	
	public void starteBestellung(String bestellNummer) {
		if (bestellNummer.startsWith("BF")) {
			bestelleingang.setBestellNr(bestellNummer);
			bestelleingang.starteBestellung();
		} else {
			int bestellnummer = Integer.parseInt(bestellNummer);
			bestellannahme = new Bestellannahme(bestellnummer);
		}
	}
} 		    
 		\end{lstlisting}
 	\end{exampleblock}
\end{frame}

\frameWithImageOnly{\BILDERPFAD/fassade4.jpg}
\note{
	Ein Heimkinosystem besteht aus vielen unterschiedlichen Komponenten. Damit man nicht bei jedem Filmabend jede Komponente einzeln ein- und ausschalten muss, bauen wir eine Fassade, mit deren Hilfe bestimmte Vorgänge gesteuert werden können. 
}

\begin{frame}{Klassendiagramm Beispiel Heimkino}
    \begin{tikzpicture}[scale=0.7, transform shape]
        \umlclass[x=-4,y=0]{Client}{
        }{
        }
    
        \umlclass[x=0,y=0]{HeimkinoFassade}{
        }{
            +filmAnschauen()
        }

        \umlclass[x=4,y=0]{Beleuchtung}{
        }{
            +dimmen()
        }

        \umlclass[x=4,y=3]{DVDPlayer}{
        }{
            +einschalten\\
            +play()
        }
        
        \umlclass[x=9,y=-2]{Projektor}{
        }{
            +einschalten()\\
            +weissesTuchHerunterfahren()
        }

        \umlclass[x=3,y=-3]{SoundSystem}{
        }{
            +einschalten()\\
            +laustarkeEinstellen(level: int)
        }        

        \umluniassoc{Client}{HeimkinoFassade}
        \umluniassoc{HeimkinoFassade}{Beleuchtung}
        \umluniassoc{HeimkinoFassade}{DVDPlayer}
        \umluniassoc{HeimkinoFassade}{Projektor}
        \umluniassoc{HeimkinoFassade}{SoundSystem}
    \end{tikzpicture}
\end{frame}

\begin{frame}[fragile]{Fassade}
	\begin{exampleblock}{HeimkinoFassade}
 		\begin{lstlisting}
public class HeimkinoFassade {
    private DVDPlayer dvdPlayer;
    private Projektor projektor;
    private Beleuchtung beleuchtung;
    private SoundSystem soundSystem;
    
    public void filmAnschauen() {
        System.out.println("Filmabend starten...");
        beleuchtung.dimmen();
        projektor.einschalten();
        projektor.weissesTuchHerunterfahren();
        soundSystem.einschalten();
        soundSystem.laustarkeEinstellen(5);
        dvdPlayer.einschalten();
        dvdPlayer.play();
    }
}		    
 		\end{lstlisting}
 	\end{exampleblock}
\end{frame}

\frameWithImageOnly{\BILDERPFAD/fassade5.jpg}
\note{
	Als nächstes soll ein Bankensystem implementiert werden. Wenn ein Kunde Geld abheben möchte, läuft ein komplexerer Prozess ab, der auf mehrere Untersysteme zugreift und Validierungen erfordert. All das wird nun in der Fassade gekapselt.
}

\begin{frame}{Klassendiagramm Beispiel Bankensystem}
    \begin{tikzpicture}[scale=0.7, transform shape]
        \umlclass[x=-3,y=3]{Client}{
        }{
        }
    
        \umlclass[x=-3,y=0]{BankFassade}{
        }{
            geldAbheben(kontonummer: int, pin: String, betrag: double)
        }

        \umlclass[x=4,y=-3]{KontoService}{
        }{
            +istKontoAktiv(kontonummer: int): boolean\\
            +hatAusreichendSaldo(betrag: double): boolean\\
            +abheben(betrag: double)
        }

        \umlclass[x=4,y=3]{KreditService}{
        }{
            +istInnerhalbDesLimits(betrag: double): boolean
        }
        
        \umlclass[x=-4,y=-3]{SicherheitsService}{
        }{
            istPinKorrekt(eingegebenePin: String): boolean
        }

        \umlclass[x=6,y=0]{TransaktionsService}{
        }{
            +durchfuehren()
        }        

        \umluniassoc{Client}{BankFassade}
        \umluniassoc{BankFassade}{KontoService}
        \umluniassoc{BankFassade}{KreditService}
        \umluniassoc{BankFassade}{SicherheitsService}
        \umluniassoc{BankFassade}{TransaktionsService}
    \end{tikzpicture}
\end{frame}

\begin{frame}[fragile]{Fassade}
	\begin{exampleblock}{BankFassade}
 		\begin{lstlisting}
public class BankFassade {
    private KontoService kontoService;
    private KreditService kreditService;
    private SicherheitsService sicherheitsService;
    private TransaktionsService transaktionsService;
    
    public void geldAbheben(int kontonummer, String pin, 
            double betrag) {
        if (!kontoService.istKontoAktiv(kontonummer)) {
            System.out.println("Konto ist nicht aktiv.");
            return;
        }
        // ...
        
        transaktionsService.durchfuehren();
        kontoService.abheben(betrag);
    }
}    
 		\end{lstlisting}
 	\end{exampleblock}
\end{frame}

\frameWithImageOnly{\BILDERPFAD/fassade6.jpg}
\note{
	Im Beispiel sollen mehrere Bibliotheken zum Auslesen von Ontologien simuliert werden: OWLApi und
	Jena. Dabei soll das resultierende System unabhängig von den entsprechenden Bibliotheken sein, so dass man beide in Kombination verwenden kann, nur eine verwenden oder neue Bibliotheken einfügen kann. Die Fassade soll die Fremd-Bibliotheken kapseln. 
}

\begin{frame}{Klassendiagramm Beispiel Ontologien}
    \begin{tikzpicture}[scale=0.7, transform shape]
        \umlclass[x=-2,y=3]{MyOntologyClass}{
        }{
        }
    
        \umlclass[x=-3,y=0]{Client}{
        }{
        }
    
        \umlclass[x=0,y=0]{OntologyFacade}{
        }{
        }

        \umlclass[x=4,y=3]{OwlOntology}{
        }{
        }

        \umlclass[x=4,y=-3]{Ontology}{
        }{
        }
        
        \umlclass[x=4,y=0]{OwlClass}{}{
        }{
        }

        \umlclass[x=1,y=-4]{OntologyClass}{}{
        }{
        }        

        \umluniassoc{Client}{OntologyFacade}
        \umluniassoc{Client}{MyOntologyClass}
        \umluniassoc{OntologyFacade}{MyOntologyClass}
        
        \umluniassoc{OntologyFacade}{OwlOntology}
        \umluniassoc{OntologyFacade}{Ontology}
        \umluniassoc{OntologyFacade}{OwlClass}
        \umluniassoc{OntologyFacade}{OntologyClass}
    \end{tikzpicture}
\end{frame}
\note{
	Im Beispiel sollen mehrere Bibliotheken zum Auslesen von Ontologien simuliert werden: OWLApi und
	Jena. Beide enthalten unterschiedliche Klassen und haben unterschiedliche Methoden. Die beiden
	Systeme sollen hinter einer Fassade versteckt werden. Zusätzlich soll die Fassade die Ergebnistypen
	in neue Klassen überführen, die unabhängig vom verwendeten Framework sind. So kann der Client
	unabhängig gehalten werden vom konkreten Framework, die Frameworks werden austauschbar und es ist 
	möglich, dass der Client für einige Anfragen die OWLApi und für andere Jena verwendet. So können
	die Stärken des jeweiligen Frameworks genutzt werden, ohne dass Änderungen am Client vorgenommen
	werden müssen.
}

\begin{frame}[fragile]{Jena}
	\begin{exampleblock}{Ontology}
 		\begin{lstlisting}
public List<OntologyClass> readAllOntologyClassess() {
    System.out.println("Ich lese alle Ontologien mit Jena");
    
    List<OntologyClass> ontologyClasses = new ArrayList<>();
    OntologyClass ontologyClass1 = new OntologyClass(4);
    OntologyClass ontologyClass2 = new OntologyClass(5);
    ontologyClasses.add(ontologyClass1);
    ontologyClasses.add(ontologyClass2);

    System.out.println("Ontologie 1: " + ontologyClass1);
    System.out.println("Ontologie 2: " + ontologyClass2);
    
    return ontologyClasses;
} 		    
 		\end{lstlisting}
 	\end{exampleblock}
\end{frame}

\begin{frame}[fragile]{OwlApi}
	\begin{exampleblock}{OwlOntology}
 		\begin{lstlisting}
public List<OwlClass> loadAllOntologyClasses() {
    uuid = UUID.randomUUID();
    System.out.println("Ich lese alle Ontologien mit OWL API");
    
    List<OwlClass> ontologyClasses = new ArrayList<>();
    OwlClass owlClass1 = new OwlClass();
    OwlClass owlClass2 = new OwlClass();
    ontologyClasses.add(owlClass1);
    ontologyClasses.add(owlClass2);

    System.out.println("Ontologie 1: " + owlClass1);
    System.out.println("Ontologie 2: " + owlClass2);
    
    return ontologyClasses;
} 		    
 		\end{lstlisting}
 	\end{exampleblock}
\end{frame}

\begin{frame}[fragile]{Fassade}
	\begin{exampleblock}{OntologyFassade}
 		\begin{lstlisting}
public List<MyOntologyClass> readAllOntologies() {
    return readAllOntologiesFromJena();
}

private List<MyOntologyClass> readAllOntologiesFromOwlApi() {
    List<OwlClass> owlClasses = owlOntology.
        loadAllOntologyClasses();
    List<MyOntologyClass> myOntologyClasses = new ArrayList<>();
    
    owlClasses.forEach(owlClass -> { 
        final String iri = owlClass.getIri().toString();
        myOntologyClasses.add(new MyOntologyClass(iri));
    });
    
    return myOntologyClasses;
} 		    
 		\end{lstlisting}
 	\end{exampleblock}
\end{frame}

